\par\Needspace{6\baselineskip}\medskip
\begin{center}\hypertarget{native-198-article148}{}\pdfbookmark[2]{148.}{article-148}148.\end{center}
Sind die aus drei Beobachtungen abzuleitenden Bahnelemente noch völlig unbekannt (für diesen Fall eignet sich unser Verfahren besonders), so werden in der ersten Hypothese, wie bereits bemerkt, für $P$ und $Q$ die Näherungswerte $\dfrac{\theta''}{\theta}$ und $\theta\theta''$ genommen. Dabei sind $\theta$ und $\theta''$ vorläufig aus den unkorrigierten Zeitintervallen abzuleiten. Drückt man deren Verhältnisse zu den korrigierten Intervallen durch $\mu:1$ beziehungsweise $\mu'':1$ aus, so erhält man in der ersten Hypothese
\[
\begin{aligned}
X&=\log\mu-\log\mu''+\log\eta-\log\eta''\\
Y&=-\log\mu-\log\mu''-\log\eta-\log\eta''\\
&\quad+\mathrm{Comp.}\log\cos f+\mathrm{Comp.}\log\cos f'\\
&\quad+\mathrm{Comp.}\log\cos f''+2\log r'-\log r-\log r''.
\end{aligned}
\]
Die Logarithmen der Größen $\mu,\mu''$ fallen gegenüber den übrigen Termen nicht ins Gewicht. Die beiden positiven Größen $\log\eta$ und $\log\eta''$ heben sich in $X$ teilweise gegenseitig auf, besonders wenn die Zeitintervalle nahezu gleich sind. Daher nimmt $X$ einen kleinen, bald positiven, bald negativen Wert an. Bei $Y$ gleichen die negativen Terme $\log\eta$ und $\log\eta''$ die positiven Terme $\mathrm{Comp.}\log\cos f$, $\mathrm{Comp.}\log\cos f'$, $\mathrm{Comp.}\log\cos f''$ zwar zum Teil aus, doch weniger vollständig, da die positiven meist deutlich überwiegen. Über das Vorzeichen von $\log\dfrac{r'r'}{rr''}$ lässt sich im Allgemeinen nichts sagen.

Ist nun die heliozentrische Bewegung zwischen den Beobachtungen mäßig, so braucht man nur selten bis zur vierten Hypothese zu gehen. Meist liefert die dritte, häufig
\SourceLeaf{199}schon die zweite eine hinreichende Genauigkeit, und bisweilen kann man sich sogar mit den Zahlen aus der ersten Hypothese begnügen. Es ist stets zweckmäßig, darauf zu achten, welchen Grad an Genauigkeit die Beobachtungen besitzen. Es wäre eine undankbare Mühe, in der Rechnung eine hundert- oder tausendmal höhere Genauigkeit anzustreben, als die Beobachtungen zulassen. In solchen Dingen schärft häufige praktische Übung das Urteil besser als Regeln, und wer geübt ist, erwirbt leicht ein Gespür dafür, wann man die Rechnung beenden sollte.
